% !TeX program = xelatex
% UTF-8. Complete transcription of the eight handwritten source pages.
% Standalone XeLaTeX/XePersian document; also usable with \input after xepersian.
% pic01.jpg = 20261013_024041820.jpg
% pic02.jpg = 20261013_024049110.jpg
% pic03.jpg = 20261013_024053464.jpg
% pic04.jpg = 20261013_024059728.jpg
% pic05.jpg = 20261013_024104790.jpg
% pic06.jpg = 20261013_024111483.jpg
% pic07.jpg = 20261013_024117391.jpg
% pic08.jpg = 20261013_024124243.jpg
% Medical statements and terminology are transcribed, not revised.
\ifdefined\settextfont
  \def\SleepStandalone{0}
\else
  \def\SleepStandalone{1}
  \documentclass[11pt,a4paper]{article}
  \usepackage[margin=2cm]{geometry}
  \usepackage{amsmath,amssymb,enumitem}
  \usepackage{xepersian}
  \IfFontExistsTF{Amiri}{\settextfont{Amiri}}{\settextfont{FreeSerif}}
  \setlatintextfont{FreeSerif}
  \setlist[itemize]{itemsep=3pt,topsep=4pt,parsep=0pt}
\fi
\providecommand{\en}[1]{\lr{#1}}
\newcommand{\sleeparr}{\ensuremath{\leftarrow}}
\newcommand{\sleeppage}[2]{%
  \clearpage\section*{کاپلان ـ ۲۳ ـ اختلالات خواب — صفحهٔ #1}
  {\footnotesize\noindent موضوع: کاپلان ـ ۲۳ ـ اختلالات خواب\par
  \noindent تاریخ: \en{\quad / \quad /}\hfill\en{#2}\par}\medskip}
\ifnum\SleepStandalone=1 \begin{document}\fi

\sleeppage{۱}{pic01.jpg}
خواب یک پروسهٔ \en{passive} نیست بلکه با فعالیت و متابولیسم مغز همراه است.

\subsection*{اختلال بی‌خوابی — \en{Insomnia Disorder}}
\begin{itemize}
  \item عدم رضایت از کمیت و کیفیت خواب.
  \item بی‌خوابی در گذشته یک علامت مربوط به اختلالات دیگر در نظر گرفته می‌شد ولی الان یک اختلال مجزا است، حتی با \en{Comorbidity} باشد.
  \item کرایتریا (حداقل یک مورد):
  \begin{itemize}
    \item دشواری در به خواب رفتن.
    \item دشواری در \en{maintaining sleep}.
    \item زود بیدار شدن صبح و عدم توانایی دوباره خوابیدن.
  \end{itemize}
  \item در کودکان می‌تواند به صورت مقاومت در برابر خوابیدن در زمان مورد نظر مراقب کودک باشد.
  \item حداقل ۳ روز در هفته، به مدت حداقل ۳ ماه.
  \item نباید بهتر توسط یک اختلال دیگر توجیه شود.
  \item اگر فرصت مناسب خواب حاصل شد بیمار نباید بتواند بخوابد.
  \item پس از تأیید تشخیص، پزشک باید موارد زیر را نیز مشخص کند:
  \begin{itemize}
    \item اختلالات هم‌زمان.
    \item اپیزودیک بودن.
    \item ریکارنس (تعداد اپیزودها در ۱ سال).
  \end{itemize}
\end{itemize}

\sleeppage{۲}{pic02.jpg}
{\footnotesize\noindent یادداشت گوشهٔ بالای صفحه: ۶۹۸۱\par}
\subsection*{\en{subtype}های اختلال بی‌خوابی (\en{Insomnias})}
\begin{itemize}
  \item بی‌خوابی به علت دارو یا مواد مخدر.
  \item بی‌خوابی کوموربید با بیماری‌های طبی.
  \item بی‌خوابی کوموربید با بیماری‌های روان‌پزشکی.
  \item \en{Behavioral Insomnia of childhood}.
  \item بهداشت خواب ناکافی.
  \item بی‌خوابی پارادوکسیکال \sleeparr{} انکار.
  \item بی‌خوابی ایدیوپاتیک.
  \item بی‌خوابی سایکوفیزیولوژیکال.
\end{itemize}
\subsection*{\en{W/U} بیمار با اختلال بی‌خوابی}
\begin{itemize}
  \item \en{TFT, HbA1c, CBC, Serum Iron, Ferritin, LFT, BUN, Cr}.
  \item تهیهٔ \en{Sleep Diaries} \sleeparr{} \en{ISI, Epworth}.
\end{itemize}
\subsection*{درمان \en{Insomnia}}
\begin{itemize}
  \item قدم اول درمان، غیردارویی است.
  \item در صورت شروع دارو در اولین فرصت به صورت تدریجی قطع شود.
  \item درمان دارویی:
\end{itemize}

\sleeppage{۳}{pic03.jpg}
\subsection*{درمان دارویی — ادامه}
\begin{itemize}
  \item \en{BZRAs} دارو سداتیو ـ هایپنوتیک استاندارد برای درمان بی‌خوابی هستند.
  \begin{itemize}
    \item ارجح نسبت به \en{BZDs} / \textit{[خوانش نامطمئن: قابل تداوم تجویز کنیم]}.
  \end{itemize}
  \item خط اول:
  \begin{itemize}
    \item آنتی‌هیستامین \en{low dose} \sleeparr{} \en{Doxepin}.
    \item آگونیست رسپتور ملاتونین \sleeparr{} \en{Ramelteon}.
    \item آنتاگونیست اورکسین \sleeparr{} \en{Suvorexant}.
  \end{itemize}
  \item ضدافسردگی‌های سداتیو \sleeparr{} توصیه نمی‌شوند.
  \item آنتی‌سایکوتیک‌های نسل ۲ \sleeparr{} اولانزاپین، ریسپریدون، کوئتیاپین.
  \begin{itemize}
    \item خط اول در ساب‌تایپ پارادوکسیکال با دلیوژنال فیچر.
    \item در سایر موارد خط دوم.
  \end{itemize}
  \item ملاتونین \sleeparr{} برای تنظیم ریتم سیرکادین.
  \item تمام موارد صرفاً برای درمان کوتاه‌مدت \en{FDA} دارند، بجز:
  \begin{itemize}
    \item زولپیدم، ازوپیکلون، راملتئون.
  \end{itemize}
\end{itemize}
\subsection*{اختلال پرخوابی — \en{Hypersomnolence disorder}}
\begin{itemize}
  \item یک اختلال نسبتاً نادر است.
  \item کرایتریای تشخیص \sleeparr{} باید همهٔ موارد باشند.
\end{itemize}
\begin{enumerate}
  \item تأخیر در بیداری (حداقل یک مورد):
  \begin{enumerate}[label=\en{\alph*)}]
    \item اپیزودهای مکرر خواب در طول یک روز.
    \item اپیزود طولانی (\en{$>9$} ساعت) خواب در یک روز که رفرشینگ نیست.
    \item دشواری در به دست آوردن وضعیت بیداری کامل.
  \end{enumerate}
  \item حداقل ۳ بار در هفته برای حداقل ۳ ماه.
\end{enumerate}

\sleeppage{۴}{pic04.jpg}
\subsection*{کرایتریای اختلال پرخوابی — ادامه}
\begin{enumerate}[start=3]
  \item باید همراه با تأثیر شدید بر بیداری، شناخت، عملکرد شغلی، اجتماعی باشد.
  \item نباید قابل ارتباط به یک اختلال دیگر باشد.
  \item نباید قابل ارتباط به دارو یا مواد مخدر باشد.
\end{enumerate}
\begin{itemize}
  \item پس از تأیید تشخیص باید کوموربیدیتی‌ها را مشخص کرد.
  \item طول مدت:
  \begin{itemize}
    \item حاد: کمتر از ۱ ماه.
    \item تحت‌حاد: ۱ تا ۳ ماه.
    \item پرسیستنت: بیشتر از ۳ ماه.
  \end{itemize}
  \item شدت:
  \begin{itemize}
    \item خفیف: ۱ تا ۲ روز در هفته.
    \item متوسط: ۳ تا ۴ روز در هفته.
    \item شدید: ۵ تا ۷ روز در هفته.
  \end{itemize}
  \item ممکن است به دنبال یک عفونت ویروسی شروع شود: \en{IMN}، پنومونی ویروسی، \en{GBS}.
  \item فرم خانوادگی با الگوی وراثت \en{AD} نیز دارد.
\end{itemize}
\subsection*{سندرم \en{Klein-Levin}}
\begin{itemize}
  \item اغلب در پسران نوجوان.
  \item اپیزودهای خواب ۱۸ تا ۲۰ ساعته.
  \item اشتها زیاد، هایپرسکشوالیتی و پرخاشگری (دیس‌این‌هیبیشن).
  \item اپیزودها بین چند روز تا چند هفته طول می‌کشند.
  \item تا ۱۰ اپیزود در سال.
\end{itemize}

\sleeppage{۵}{pic05.jpg}
\begin{itemize}
  \item اختلالات طبی که پرخوابی ایجاد می‌کنند \sleeparr{} تروما به سر، \en{CVA}، انسفالیت، پارکینسون، تومورها، بیماری‌های نورودژنراتیو.
\end{itemize}
\subsection*{پرخوابی در همراهی با منس}
\begin{itemize}
  \item یک هفته طول می‌کشد. با شروع منس برطرف می‌شود.
  \item درمان با \en{OCP} مؤثر است.
  \item ثانویه به \en{hormone Imbalance} است.
\end{itemize}
\subsection*{درمان اختلال پرخوابی}
\begin{itemize}
  \item درمان معمولاً دارویی است.
  \item محرک کلاسیک \sleeparr{} آمفتامین‌ها، متیل‌فنیدیت‌ها.
  \item داروهای \en{wake-promoting} \sleeparr{} مودافینیل، \en{pitolisant, solriamfetol}.
\end{itemize}
\subsection*{نارکولپسی}
تتراد نارکولپسی:
\begin{itemize}
  \item \en{Excessive Sleepiness} — علائم کلیدی.
  \item \en{Cataplexy} — علائم کلیدی.
  \item \en{Sleep Paralysis}.
  \item \en{Hypnagogic Hallucinations}.
\end{itemize}
\begin{itemize}
  \item کاتاپلکسی: از دست رفتن ناگهانی تون عضلانی.
  \begin{itemize}
    \item با هوشیاری کامل بیمار همراه است.
  \end{itemize}
\end{itemize}

\sleeppage{۶}{pic06.jpg}
\subsection*{یادداشت‌های بالای صفحه}
\begin{itemize}
  \item توهم هیپناگوژیک \sleeparr{} هنگام خواب رفتن.
  \item هیپنوپومپیک \sleeparr{} هنگام بیدار شدن.
\end{itemize}
\subsection*{کاتاپلکسی — ادامه}
\begin{itemize}
  \item به طور شایع به دنبال خنده یا احساسات شدید.
  \item شدت از ضعف خفیف تا کوادری‌پلژی.
  \item چند ثانیه تا چند دقیقه.
  \item بیمار نمی‌تواند صحبت کند.
\end{itemize}
\subsection*{یافته‌های پلی‌سومنوگرافی}
\begin{itemize}
  \item \en{REM} بلافاصله پس از به خواب رفتن رخ می‌دهد.
  \begin{itemize}
    \item چه در خواب روز، چه شب.
  \end{itemize}
\end{itemize}
خواب شبان در نارکولپسی معمولاً منقطع است.

بیمار نارکولپسی معمولاً به دپرشن نیز رنج می‌برد.

می‌تواند یک اختلال خودایمنی باشد که در آن نورون‌های تولیدکنندهٔ هایپوکرتین از بین می‌روند.

\subsection*{انواع نارکولپسی براساس \en{DSM-5-TR}}
\begin{latin}
\begin{itemize}
  \item Narcolepsy w/o Cataplexy but with hypocretin deficiency.
  \item Narcolepsy with Cataplexy but w/o hypocretin deficiency, AD; cerebellar Ataxia, deafness.
  \item AD Narcolepsy, obesity, DM type 2.
  \item Narcolepsy secondary to another medical disorder.
\end{itemize}
\end{latin}

\sleeppage{۷}{pic07.jpg}
\subsection*{کرایتریای تشخیص نارکولپسی}
\begin{itemize}
  \item \en{A}: اپیزودهای تکرارشونده از نیاز غیرقابل جلوگیری به خواب، به خواب رفتن، چرت زدن که در طول روز تکرار شود.
  \begin{itemize}
    \item حداقل ۳ روز در هفته در ۳ ماه گذشته.
  \end{itemize}
  \item \en{B}: وجود حداقل یک مورد:
  \begin{itemize}
    \item اپیزودهای کاتاپلکسی.
    \item هایپوکرتین دفیشنسی در \en{CSF} \sleeparr{} \en{$\leq\frac{1}{3}\,\mathrm{NL}$}، \en{$\leq110\,\mathrm{pg/ml}$}؛ که در زمینهٔ تروما، التهاب یا عفونت نباشد.
    \item در پلی‌سومنوگرافی: \en{REM} ظرف ۱۵ دقیقه از شروع خواب.
  \end{itemize}
\end{itemize}
\subsection*{نارکولپسی در زمینهٔ سایر بیماری‌ها}
معمولاً همراه با کاتاپلکسی است.
\begin{itemize}
  \item تومور و سارکوئیدوز هیپوتالاموس.
  \item \en{MS} هیپوتالاموس.
  \item سندرم پارانئوپلاستیک با آنتی‌بادی \en{anti-Ma2}.
  \item بیماری \en{Niemann-Pick} تایپ \en{C} (خصوصاً در کودکان زیر ۵ سال).
  \item سندرم \en{Coffin-Lowry}.
\end{itemize}
\subsection*{بیماری‌هایی که همراه بدون کاتاپلکسی است}
\begin{itemize}
  \item پارکینسون، سندرم پرادر ـ ویلی، دیستروفی مایوتونیک، \en{MS}.
\end{itemize}

\sleeppage{۸}{pic08.jpg}
\subsection*{درمان نارکولپسی}
\begin{itemize}
  \item خط اول: مودافینیل.
  \item محرک‌ها \sleeparr{} در صورت شکست مودافینیل \sleeparr{} خط دوم.
  \item درمان‌های جدید خط دوم:
  \begin{itemize}
    \item \en{Pitolisant} (آگونیست معکوس انتخابی گیرندهٔ هیستامین \en{$H_3$}).
    \item \en{Solriamfetol} (مهارکنندهٔ انتخابی بازجذب \en{NEP} و دوپامین).
  \end{itemize}
  \item برای درمان کاتاپلکسی:
  \begin{itemize}
    \item سدیم اکسیبات (\en{Oxybate}).
    \item داروهای سرکوب‌کنندهٔ \en{REM} (بسیاری از ضدافسردگی‌ها).
    \begin{itemize}
      \item \en{SSRIs}.
    \end{itemize}
  \end{itemize}
\end{itemize}
\subsection*{اختلالات تنفسی مرتبط با خواب}

% The source ends at the above heading. Dates are blank on all eight pages.
\ifnum\SleepStandalone=1 \expandafter\enddocument\fi
